(1) 对并联节点应用 KCL：
\[
C\dot v+\frac vR+i_L=0,\qquad L\dot i_L=v.
\]
对第一式求导并消去电感电流，得到
\[
\ddot v+\frac1{RC}\dot v+\frac1{LC}v=0,\qquad
\dot v(0)=-\frac{V_0/R+I_0}{C}.
\]
不能将并联电路直接误写为串联 RLC 的阻尼系数 \(R/L\)。

(2) 特征根
\[
r_\pm=-\frac1{2RC}\pm\sqrt{\frac1{4R^2C^2}-\frac1{LC}}.
\]
若 \(r_+\ne r_-\)，则 \(v(t)=A_+e^{r_+t}+A_-e^{r_-t}\)，其中
\[
A_+=\frac{\dot v(0)-r_-V_0}{r_+-r_-},\qquad
A_-=\frac{r_+V_0-\dot v(0)}{r_+-r_-}.
\]
判别式小于零是欠阻尼，大于零是过阻尼；等于零时根 \(r=-1/(2RC)\) 且 \(v(t)=[V_0+(\dot v(0)-rV_0)t]e^{rt}\)。所有情形实部为负，故 \(v,i_L\to0\)。

(3) 将 KCL 乘以 \(v\)，结合 \(L i_L\dot i_L=i_Lv\)，得
\[
\frac d{dt}\left(\frac12Cv^2+\frac12Li_L^2\right)=-\frac{v^2}{R}\le0.
\]
因终态两种储能均为零，积分获得
\[
\int_0^\infty\frac{v(t)^2}{R}\,dt=\frac12C V_0^2+\frac12L I_0^2.
\]
这是总能量守恒与电阻耗散的精确结果，而非只在弱阻尼下成立。

(4) 对电感方程积分：\(L[i_L(\infty)-I_0]=\int_0^\infty v(t)\,dt\)，因此
\[
\int_0^\infty v(t)\,dt=-L I_0.
\]
如果 \(I_0=0\)，积分为零，即便 \(V_0\ne0\)，因为衰减过程中的电压会跨过零并使正负面积抵消。该积分有符号，不等于正定的总耗散功。
